% ---
% Subsección 2.1.6 de los Antecedentes (nueva en TT2; antes 2.1.7 en el borrador del 2026-10-14).
% Entradas de la bitácora relacionadas: D-029 (barrido del factor de escala), D-032 (protocolo de la
% geometric difference); H-031 (kernel concentrado), H-038, H-048 (etiquetas construidas).
% Cada atribución lleva su respaldo verificable con Code/tools/verificar_citas_fichas.py.
% ---

\subsubsection{Límites de la ventaja cuántica con datos clásicos}
\label{sec:limites_ventaja}

La búsqueda de una aceleración cuántica (\textit{quantum speed-up}) en aprendizaje supervisado ha producido demostraciones formales, pero restringidas a familias de problemas con una estructura matemática muy particular. Conocer esas demostraciones y sus límites es necesario para valorar qué cabe esperar del aprendizaje automático cuántico sobre datos clásicos genéricos, como las características radiómicas extraídas de mamografías.

\textbf{Ventaja cuántica a partir de una estructura algebraica en la tarea}

Las demostraciones más rigurosas de ventaja cuántica en aprendizaje supervisado con datos clásicos provienen de tareas construidas para que la regla de clasificación esconda una estructura matemática.
% respaldo: «In this paper, we establish a rigorous quantum speed-up for supervised classification using a general-purpose quantum learning algorithm that only requires classical access to data.» (Liu2021Speedup, p. 1)
Liu et al. demostraron una aceleración cuántica rigurosa para clasificación supervisada con un algoritmo que solo requiere acceso clásico a los datos \cite{Liu2021Speedup}.
% respaldo: «Our quantum classifier is a conventional support vector machine that uses a fault-tolerant quantum computer to estimate a kernel function.» (Liu2021Speedup, p. 1)
Su clasificador es una máquina de vectores de soporte convencional que usa una computadora cuántica tolerante a fallos para estimar la función kernel \cite{Liu2021Speedup}.
% respaldo: «We construct a family of datasets and show that no classical learner can classify the data inverse-polynomially better than random guessing, assuming the widely-believed hardness of the discrete logarithm problem.» (Liu2021Speedup, p. 1)
Bajo el supuesto de que el problema del logaritmo discreto es difícil para las computadoras clásicas, construyeron una familia de conjuntos de datos en la que ningún clasificador clásico eficiente supera al azar por más que un margen inverso-polinomial \cite{Liu2021Speedup}.
% respaldo: «the training data samples are clustered in two intervals: [s, s + p−3 2 ] with label +1, and [s + p−1 2 , s + p −2] with label −1, for the unknown s ∈Z∗ p (which defines the unknown concept fs).» (Liu2021Speedup, p. 11)
La dificultad está en la estructura algebraica de la tarea: vistos a través del logaritmo discreto, los datos de cada clase forman dos intervalos \cite{Liu2021Speedup}.
% respaldo: «While the learning problem we have presented here that demonstrates an exponential quantum speed-up is not practically motivated, our result sets a positive theoretical foundation for the search of practical quantum advantage in machine learning.» (Liu2021Speedup, p. 7)
Los propios autores aclaran que su problema no tiene motivación práctica y que su resultado es un fundamento teórico para buscar ventajas prácticas \cite{Liu2021Speedup}.

% respaldo: «It is an important challenge to find quantum kernels that provide an advantage in the classification of real-world data.» (Glick2024Covariant, p. 1)
Glick et al. parten de que encontrar kernels cuánticos que den una ventaja en la clasificación de datos reales es un reto importante \cite{Glick2024Covariant}.
% respaldo: «We introduce a class of quantum kernels that can be used for data with a group structure.» (Glick2024Covariant, p. 1)
Para avanzar en esa dirección introducen una clase de kernels cuánticos, llamados covariantes, para datos con estructura de grupo \cite{Glick2024Covariant}.
% respaldo: «The learning problem we have studied experimentally serves as an abstract representation of many group theoretic problems, i.e. classification of data that is close to cosets of a group.» (Glick2024Covariant, p. 4)
El problema que evalúan experimentalmente es la clasificación de datos cercanos a las clases laterales (\textit{cosets}) de un grupo \cite{Glick2024Covariant}.
% respaldo: «We implement the learning algorithm with 27 qubits on a superconducting processor.» (Glick2024Covariant, p. 1)
% respaldo: «The considered circuit, although supported on a larger number of qubits, is still sufficiently shallow so that classical simulation methods can be applied [37].» (Glick2024Covariant, p. 5)
Lo implementaron con 27 qubits en un procesador superconductor, y señalan que el circuito es lo bastante poco profundo como para que puedan aplicarse métodos de simulación clásica \cite{Glick2024Covariant}.

\textbf{Ventaja cuántica con datos que ya son estados cuánticos}

Una segunda clase de resultados con ventaja exponencial corresponde a escenarios en los que las entradas no son datos clásicos, sino estados de un sistema cuántico.
% respaldo: «An experimental setup that transduces data from a physical system to a stable quantum memory, and processes that data using a quantum computer, could have significant advantages over conventional experiments in which the physical system is measured and the outcomes are processed using a classical computer.» (Huang2022Experiments, p. 1)
Huang et al. muestran que un experimento que transfiere los datos de un sistema físico a una memoria cuántica estable y los procesa en una computadora cuántica puede tener ventajas significativas sobre uno en el que el sistema se mide y los resultados se procesan en una computadora clásica \cite{Huang2022Experiments}.
% respaldo: «The exponential advantage holds in predicting properties of physical systems, performing quantum principal component analysis on noisy states, and learning approximate models of physical dynamics.» (Huang2022Experiments, p. 1)
La ventaja exponencial se da en la predicción de propiedades de sistemas físicos, en el análisis cuántico de componentes principales sobre estados ruidosos y en el aprendizaje de modelos aproximados de dinámicas físicas \cite{Huang2022Experiments}.
% respaldo: «showing that a sufficiently large quantum memory is needed to achieve this task in the quantum-enhanced scenario.» (Huang2022Experiments, p. 3)
Para alcanzarla hace falta una memoria cuántica suficientemente grande \cite{Huang2022Experiments}. Este resultado no se extiende a los datos del presente trabajo: una mamografía digitalizada y los descriptores radiómicos que se extraen de ella son registros clásicos de una medición ya realizada, no estados cuánticos que puedan cargarse en una memoria cuántica.

\textbf{Datos clásicos codificados en ángulos}

Para datos clásicos arbitrarios codificados mediante rotaciones, la literatura reciente describe limitaciones que condicionan la posibilidad de una ventaja.
% respaldo: «with the help of data, classical machine learning models can be competitive with quantum models even if they are tailored to quantum problems.» (huang2021power, p. 1)
Huang et al. muestran que, con ayuda de datos de entrenamiento, los modelos clásicos pueden competir con los cuánticos incluso en problemas diseñados para estos últimos \cite{huang2021power}.
% respaldo: «we find that a variety of common quantum models in the literature perform similarly or worse than classical ML on both classical and quantum datasets due to a small geometric difference.» (huang2021power, p. 2)
% respaldo: «The small geometric difference is a consequence of the exponentially large Hilbert space employed by existing quantum models, where all inputs are too far apart.» (huang2021power, p. 2)
En sus experimentos numéricos, varios modelos cuánticos habituales rinden igual o peor que los clásicos por tener una diferencia geométrica $g$ pequeña, que atribuyen a que en un espacio de Hilbert exponencialmente grande todas las entradas quedan demasiado lejos entre sí \cite{huang2021power}.
% respaldo: «Only when the geometric difference g is large (projected kernel method with embedding E3) can we see some mild advantage over the best classical method.» (huang2021power, p. 7)
Solo cuando $g$ fue grande, con un kernel proyectado, observaron una ventaja leve sobre el mejor método clásico \cite{huang2021power}.

% respaldo: «We show that, under certain conditions, values of quantum kernels over different input data can be exponentially concentrated (in the number of qubits) towards some fixed value.» (Thanasilp2024, p. 1)
Thanasilp et al. demostraron que, bajo ciertas condiciones, los valores de un kernel cuántico sobre entradas distintas se concentran exponencialmente, con el número de qubits, en torno a un valor fijo \cite{Thanasilp2024}.
% respaldo: «We identify four sources that can lead to concentration including expressivity of data embedding, global measurements, entanglement and noise.» (Thanasilp2024, p. 1)
Identifican cuatro fuentes de esta concentración: la expresividad del \textit{embedding}, las mediciones globales, el entrelazamiento y el ruido \cite{Thanasilp2024}.
% respaldo: «In such cases, the number of shots required to resolve the kernels to a sufficiently high accuracy scales exponentially.» (Thanasilp2024, p. 2)
En esos casos, el número de mediciones (\textit{shots}) necesario para resolver los valores del kernel con suficiente precisión crece exponencialmente \cite{Thanasilp2024}.
% respaldo: «Thus on training with a polynomial number of measurements, one ends up with a trivial model where the predictions on unseen inputs are independent of the input data.» (Thanasilp2024, p. 1)
Con un número polinomial de mediciones, el modelo entrenado resulta trivial: sus predicciones sobre entradas nuevas no dependen de esas entradas \cite{Thanasilp2024}.

% respaldo: «However, having an exponentially large feature space renders the problem of generalization hard.» (Kubler2021Inductive, p. 1)
Desde el análisis del sesgo inductivo, Kübler et al. señalan que un espacio de características exponencialmente grande vuelve difícil la generalización \cite{Kubler2021Inductive}.
% respaldo: «Furthermore, being able to evaluate inner products in high dimensional spaces efficiently by itself does not guarantee a quantum advantage» (Kubler2021Inductive, p. 1)
Poder evaluar eficientemente productos internos en espacios de dimensión alta no garantiza, por sí solo, una ventaja cuántica \cite{Kubler2021Inductive}.
% respaldo: «To make learning from data feasible we cannot use the full expressive power of the quantum Hilbert space but instead only very restricted embeddings allow to learn from data.» (Kubler2021Inductive, p. 7)
Para que el aprendizaje sea viable no puede usarse toda la expresividad del espacio de Hilbert, sino solo \textit{embeddings} muy restringidos \cite{Kubler2021Inductive}.
% respaldo: «To summarize, we conclude that there is no indication that quantum machine learning will substantially improve supervised learning on classical datasets.» (Kubler2021Inductive, p. 10)
Concluyen que no hay indicios de que el aprendizaje automático cuántico vaya a mejorar sustancialmente el aprendizaje supervisado sobre conjuntos de datos clásicos \cite{Kubler2021Inductive}.

\textbf{El ancho de banda y la evidencia empírica}

Entre los dos extremos de un kernel (uno cercano a la identidad, en el que cada entrada solo se parece a sí misma, y uno casi constante, en el que todas se parecen) la literatura ha propuesto controlar el ancho de banda como hiperparámetro.
% respaldo: «In this work we introduce the hyperparameter controlling the bandwidth of a quantum kernel and show that it controls the expressivity of the resulting model.» (Shaydulin2022Bandwidth, p. 1)
Shaydulin y Wild introdujeron el hiperparámetro de ancho de banda (\textit{bandwidth}) de un kernel cuántico y mostraron que controla la expresividad del modelo \cite{Shaydulin2022Bandwidth}.
% respaldo: «We show that for the quantum feature maps considered, the bandwidth can be controlled by rescaling the datapoints.» (Shaydulin2022Bandwidth, p. 2)
% respaldo: «Large values of t lead to the kernel’s being too “narrow” (small bandwidth), making learning impossible.» (Shaydulin2022Bandwidth, p. 3)
En los mapas que estudian, el ancho de banda se controla reescalando los datos: un factor de escala grande hace el kernel demasiado estrecho, es decir, de ancho de banda pequeño, y el aprendizaje se vuelve imposible \cite{Shaydulin2022Bandwidth}.
% respaldo: «We use extensive numerical experiments with multiple quantum kernels and classical datasets to show consistent change in the model behavior from underfitting (bandwidth too large) to overfitting (bandwidth too small), with optimal generalization in between.» (Shaydulin2022Bandwidth, p. 1)
En sus experimentos numéricos, el modelo pasa del subajuste, con un ancho de banda demasiado grande, al sobreajuste, con uno demasiado pequeño, y generaliza mejor en un punto intermedio \cite{Shaydulin2022Bandwidth}.
% respaldo: «We reproduce these negative results and show that if the kernel bandwidth is optimized, the performance instead improves with growing qubit count and becomes competitive with the best classical methods.» (Shaydulin2022Bandwidth, p. 1)
Al reproducir los resultados negativos de Huang et al., observaron que con el ancho de banda optimizado el rendimiento mejora con el número de qubits y se vuelve competitivo con los mejores métodos clásicos \cite{Shaydulin2022Bandwidth}.

% respaldo: «Specifically, we show that changing the value of the bandwidth can take a model from provably not being able to generalize to any target function to good generalization for well-aligned targets.» (Canatar2023Bandwidth, p. 1)
Canatar et al. muestran que cambiar el ancho de banda puede llevar a un modelo desde no poder generalizar a ninguna función objetivo, de forma demostrable, hasta generalizar bien para objetivos alineados con el kernel \cite{Canatar2023Bandwidth}.
% respaldo: «Our analysis shows how the bandwidth controls the spectrum of the kernel integral operator and thereby the inductive bias of the model.» (Canatar2023Bandwidth, p. 1)
Su análisis explica el mecanismo: el ancho de banda controla el espectro del operador integral del kernel y, con él, el sesgo inductivo del modelo \cite{Canatar2023Bandwidth}.
% respaldo: «On the other hand, bandwidth enables sufficient decay in the spectrum of the kernel, and Eq. 9 shows that P ∼O(nl) samples yield an excess generalization error» (Canatar2023Bandwidth, p. 8)
Con un ancho de banda adecuado, el espectro decae lo suficiente como para que baste un número de muestras polinomial en el número de qubits \cite{Canatar2023Bandwidth}.

En el presente trabajo, el factor de escala que multiplica los ángulos de codificación desempeña el papel del factor de Shaydulin y Wild. Con el valor del diseño, igual a uno, el kernel cuántico queda en el extremo estrecho; reducirlo equivale a ensanchar el ancho de banda.

Las evaluaciones comparativas recientes insisten en contrastar los modelos cuánticos con métodos clásicos convencionales.
% respaldo: «on the PennyLane software framework and use it to conduct a large-scale study that systematically tests 12 popular quantum machine learning models on 6 binary classification tasks used to create 160 individual datasets.» (Bowles2024Benchmarking, p. 1)
% respaldo: «We find that overall, out-of-the-box classical machine learning models outperform the quantum classifiers.» (Bowles2024Benchmarking, p. 1)
En un estudio que evaluó 12 modelos cuánticos sobre 160 conjuntos de datos, Bowles et al. encontraron que, en conjunto, los modelos clásicos estándar, sin ajustes específicos, superan a los clasificadores cuánticos \cite{Bowles2024Benchmarking}.
% respaldo: «In all three model families, the prototypical classical models came first.» (Bowles2024Benchmarking, p. 14)
En las tres familias examinadas (redes neuronales cuánticas, métodos de kernel cuánticos y redes convolucionales cuánticas), el modelo clásico de referencia quedó en primer lugar \cite{Bowles2024Benchmarking}.
% respaldo: «Moreover, removing entanglement from a quantum model often results in as good or better performance, suggesting that “quantumness” may not be the crucial ingredient for the small learning tasks considered here.» (Bowles2024Benchmarking, p. 1)
Además, quitar el entrelazamiento de un modelo cuántico dio a menudo un rendimiento igual o mejor, lo que sugiere que lo ``cuántico'' quizá no sea el ingrediente decisivo en las tareas pequeñas que consideraron \cite{Bowles2024Benchmarking}.

% respaldo: «We argue that these challenges call for a critical debate on whether quantum advantage and the narrative of “beating” classical machine learning should continue to dominate the literature the way it does, and highlight examples for how other perspectives in existing research provide an important alternative to the focus on advantage.» (SchuldKilloran2022Advantage, p. 1)
Ante estas dificultades, Schuld y Killoran proponen debatir si la ventaja cuántica y la narrativa de ``vencer'' al aprendizaje clásico deben seguir dominando la literatura como lo hacen \cite{SchuldKilloran2022Advantage}.
% respaldo: «To state it in simple terms, quantum machine learning research is trying to beat large, high-performing algorithms for problems that are conceptually hard to study.» (SchuldKilloran2022Advantage, p. 1)
En su diagnóstico, la investigación en QML intenta superar algoritmos grandes y de alto rendimiento en problemas conceptualmente difíciles de estudiar \cite{SchuldKilloran2022Advantage}.
% respaldo: «How can we make quantum software ready for machine learning applications?.» (SchuldKilloran2022Advantage, p. 2)
Entre las preguntas alternativas que plantean está cómo preparar el software cuántico para aplicaciones de aprendizaje automático \cite{SchuldKilloran2022Advantage}.

\textbf{Conexión con el diseño del presente trabajo}

Las ventajas demostradas hasta ahora se apoyan en tareas con estructura algebraica o en entradas que ya son estados cuánticos. Para datos clásicos codificados en ángulos, la literatura revisada asocia la concentración del kernel y una diferencia geométrica pequeña con modelos clásicos que igualan o superan a los cuánticos, y ninguna de las obras revisadas da razones para anticipar una superioridad cuántica en ese régimen.

Por eso, el presente trabajo no parte de la hipótesis de una ventaja cuántica ni plantea el experimento como una competencia de rendimiento. La separabilidad de clases en los espacios inducidos por los circuitos se mide con métricas que usan las etiquetas (sección~\ref{sec:metricas}). La diferencia geométrica de Huang et al., que no usa las etiquetas, se emplea como indicador de si el kernel cuántico induce una geometría que un conjunto de kernels clásicos no reproduce. El efecto del ancho de banda se examina con un barrido del factor de escala angular.
